\section{Flood, Plague, and a Reluctant Election (589--590)}

\subsection{The eyewitness record of 590}

For the crisis that placed Gregory on the Roman throne, history possesses an
extraordinary contemporary witness written within months of the events.
Gregory of Tours opens the tenth and final book of his \work{Historia
Francorum} with the firsthand report of his own deacon, who had travelled to
Rome to obtain relics of the saints and stayed to witness both the plague and
Gregory's consecration.

According to that deacon's account, in November 589 (``the ninth month'' of the
previous year, counting from March) the Tiber burst its banks and submerged
Rome under such torrents that ancient temples collapsed and the Church's
granaries were ruined, destroying several thousand bushels of wheat. A mass of
serpents, among them a great dragon ``large as a stout beam,'' washed down the
swollen river to the sea, drowned in the surf, and drifted back upon the shore
(\work{Historia Francorum} X.1). In the middle of January 590 (the eleventh
month), the bubonic pestilence known as the \latin{lues inguinaria}---``the
plague of the groin''---broke out across the flooded city.

It struck the head of the Roman Church first. Pope Pelagius II fell ill and
died swiftly (7 or 8 February 590), fulfilling, as the chronicler remarked,
Ezekiel 9:6: ``And begin at My sanctuary.'' As the death toll mounted across
every quarter of the city, clergy and people turned with one voice to the
deacon Gregory. He knew both the streets of Rome and the palace of
Constantinople; he had governed the city as prefect and prayed over it as a
monk.

Gregory did everything in his power to escape. Fearing ``that if he assumed it,
something of the worldly pride which he had long cast away might steal back
upon him,'' he wrote privately to Emperor Maurice, imploring his imperial
friend and \latin{compater} to withhold confirmation of the election
(\work{Historia Francorum} X.1). Under sixth-century Byzantine law, no bishop
of Rome could be consecrated without the emperor's consent. Yet Germanus, the
Prefect of Rome, intercepted the courier, seized and destroyed Gregory's
letter, and dispatched only the Roman decree requesting confirmation. Maurice
rejoiced at the opportunity to elevate his former envoy and sent back the
imperial diploma commanding Gregory's consecration.

\subsection{The plague sermon and the sevenfold litany}

While Rome waited through the spring and summer of 590 for imperial
confirmation, the epidemic raged unchecked. Still a deacon and pope-elect,
Gregory stepped forward to lead the city in public penance. Gregory of Tours
preserves the full text of Gregory's urgent address to the Roman people (also
transmitted separately among Gregory's works as the \work{Oratio ad plebem de
mortalitate}, Migne PL 76, col.~1311):

\begin{quote}
Behold how all the people is smitten by the sword of the divine wrath; one
after another, they are swept away by sudden death. No gradual sickness cometh
before death, but, as ye behold, death forestalleth the slow steps of sickness.
The blow falleth; the victim is snatched away before he can turn to bewail his
sins and to repent. \dots\ Houses are left void, parents behold the funerals of
their children, and their heirs go before them to the grave. Let every one of
us therefore betake himself to lamentation and repentance before the blow is
fallen, and while time yet remaineth to weep. (\work{Historia Francorum} X.1)
\end{quote}

Yet Gregory refused to let terror collapse into despair. Reminding the Romans
that three days of repentance had spared Nineveh and a single dying prayer had
won paradise for the good thief, he summoned the entire city to assemble at
dawn on Wednesday for a \latin{litania septiformis}---a sevenfold penitential
procession. Seven distinct companies, each led by the priests of its assigned
region, were ordered to set out from seven Roman churches: the clergy from
Saints Cosmas and Damian; abbots and monks from Saints Protasius and Gervasius;
abbesses and nuns from Saints Marcellinus and Peter; children from Saints John
and Paul; laymen from Saint Stephen; widows from Saint Euphemia; and married
women from Saint Clement. All seven columns were to converge with tears and
psalmody on the basilica of Saint Mary Major (\latin{ad beatae Mariae semper
virginis genetricis domini nostri Iesu Christi basilicam}) to entreat God's
mercy (\work{Historia Francorum} X.1).

The deacon of Tours watched the three-day supplication unfold as choirs moved
through the streets chanting \latin{Kyrie eleison}. The plague did not halt at
once: ``Our deacon, who was present, declared that in the space of a single
hour, while the people were pouring forth the voice of their supplication to
the Lord, there fell dead no less than eighty persons'' (\work{Historia
Francorum} X.1). Gregory did not flinch; he kept urging the survivors not to
cease from prayer. Notably, the eyewitness from Tours---who recorded the
serpents in the Tiber, the exact assembly churches, and the eighty sudden
deaths---says nothing whatsoever of an angel sheathing a sword above Hadrian's
mausoleum. That celebrated vision belongs to a much later layer of medieval
Roman memory examined below.

\subsection{Consecration and interior exile}

When Maurice's confirmation finally reached Rome, Gregory made one last attempt
to slip away. Gregory of Tours records the scene in bare, sober strokes: ``He
was preparing flight to some place of concealment, when he was seized, carried
along, and brought into the basilica of Saint Peter, where he was consecrated
bishop, and given as pope to the city of Rome'' (\work{Historia Francorum}
X.1). The deacon of Tours hurried back from the port of Ostia just to see the
consecration with his own eyes. The date was Sunday, 3 September 590 (later
kept in the Roman Martyrology as the feast of his ordination alongside 12
March, the day of his death). Later lives by Paul the Deacon and John the
Deacon embroidered his attempted flight with a disguise, a forest cave, and a
pillar of heavenly light disclosing his hiding place (Barmby, Prolegomena); the
contemporary eyewitness knew only a reluctant deacon caught before he could
vanish.

Once consecrated, Gregory did not pretend to enjoy his elevation. In October
590, writing to Theoctista, the emperor's sister, he poured out a grief so raw
that it reads like an elegy for his lost cloister:

\begin{quote}
For I have lost the deep joys of my quiet, and seem to have risen outwardly
while inwardly falling down. Whence I grieve to find myself banished far from
the face of my Maker. \dots\ I have loved the beauty of the contemplative life
as a Rachel, barren, but keen of sight and fair, who, though in her quietude
she is less fertile, yet sees the light more keenly. But, by what judgment I
know not, Leah has been coupled with me in the night, to wit, the active life;
fruitful, but tender-eyed; seeing less, but bringing forth more. I have longed
to sit at the feet of the Lord with Mary, to take in the words of His mouth;
and lo, I am compelled to serve with Martha in external affairs, to be careful
and troubled about many things. \dots\ Lo, our most serene Lord the Emperor has
ordered an ape to be made a lion. And, indeed, in virtue of his order it can be
called a lion, but a lion it cannot be made. (\work{Ep.}\ I.5)
\end{quote}

Yet even in those first overwhelming months, Gregory set his hand firmly to the
rudder. In his formal synodical letter of February 591 (\work{Ep.}\ I.25), sent
according to ancient custom to the four Eastern patriarchs---John the Faster of
Constantinople, Eulogius of Alexandria, Gregory of Antioch, and John of
Jerusalem, pointedly adding the unjustly deposed rightful patriarch Anastasius
of Antioch---he professed the Catholic faith in words that became a benchmark
for Western orthodoxy:

\begin{quote}
I confess that I receive and revere, as the four books of the Gospel so also
the four Councils: to wit, the Nicene, in which the perverse doctrine of Arius
is overthrown; the Constantinopolitan also, in which the error of Eunomius and
Macedonius is refuted; further, the first Ephesine, in which the impiety of
Nestorius is condemned; and the Chalcedonian, in which the pravity of Eutyches
and Dioscorus is reprobated. These with full devotion I embrace, and adhere to
with most entire approval; since on them, as on a four-square stone, rises the
structure of the holy faith. (\work{Ep.}\ I.25)
\end{quote}

To the four great councils he immediately added his equal veneration of the
Fifth Ecumenical Council (Constantinople II, 553) and its condemnation of the
Three Chapters. The reluctant monk had already found his public voice.
